%%%%%%%%%%%%%%%%%%%%%%%%%%%%%%%%%%%%%%%%%%%%%%%%%%%%%%%%%%%%%%%%%%%%%%%%%%%%%%%%
% Životopis – Norbert Micheľ. Postavené na altacv.cls v1.7.3 (LianTze Lim).
% Kompilácia: ./scripts/build-project.ps1 -Project cv -MainFile norbert-michel-cv.tex
% Metadáta PDF (názov, autor, kľúčové slová): norbert-michel-cv.xmpdata
%
% Makro \fillin{...} označí text, ktorý treba ešte doplniť (v hotovom CV nesmie ostať).
%%%%%%%%%%%%%%%%%%%%%%%%%%%%%%%%%%%%%%%%%%%%%%%%%%%%%%%%%%%%%%%%%%%%%%%%%%%%%%%%
\documentclass[10pt,a4paper,ragged2e,withhyper]{altacv}

\geometry{left=1.25cm,right=1.25cm,top=1.1cm,bottom=1.0cm,columnsep=0.8cm}

\usepackage{paracol}
\usepackage[rm]{roboto}
\usepackage[defaultsans]{lato}
\renewcommand{\familydefault}{\sfdefault}
\usepackage{microtype}   % jemnejšia sadzba: menej pretečených riadkov, rovnomernejší text
\usepackage{qrcode}      % QR kód v hlavičke
\usepackage{needspace}   % položka sa nerozdelí medzi stranami

%--- Farby: pôvodná bordová paleta, ale tmavší text (lepší kontrast, aj pri tlači)
\definecolor{Ink}{HTML}{222222}
\definecolor{BodyGrey}{HTML}{474747}
\definecolor{DarkPastelRed}{HTML}{450808}
\definecolor{PastelRed}{HTML}{8F0D0D}
\definecolor{GoldenEarth}{HTML}{E7D192}
\definecolor{Gold}{HTML}{9A7A1F}
\definecolor{Sidebar}{HTML}{FAF6EC}
\colorlet{name}{Ink}
\colorlet{tagline}{PastelRed}
\colorlet{heading}{DarkPastelRed}
\colorlet{headingrule}{GoldenEarth}
\colorlet{subheading}{PastelRed}
\colorlet{accent}{PastelRed}
\colorlet{emphasis}{Ink}
\colorlet{body}{BodyGrey}

\renewcommand{\namefont}{\fontsize{26}{30}\selectfont\rmfamily\bfseries}
\renewcommand{\taglinefont}{\large\bfseries}
\renewcommand{\cvsectionfont}{\Large\rmfamily\bfseries}
\renewcommand{\cvItemMarker}{{\color{accent}\small\textbullet}}

%=== OPRAVY CHÝB Z PÔVODNEJ VERZIE ===========================================
% (1) \divider ukončí odsek -> nadpis za čiarou sa už "neprilepí" k čiare
%     (to spôsobovalo "Full-", "Web", "Ap-" trčiace na okraji strany 2)
%     + čiara je "leaders" (lepidlo), takže pri zlome strany zmizne a nezostane osamotená na spodku
\renewcommand{\divider}{\par\needspace{8\baselineskip}\vskip6pt
  \cleaders\vbox to 0.6pt{\vss\hbox to\linewidth{\textcolor{body!25}{\hdashrule{\linewidth}{0.6pt}{0.5ex}}}}\vskip0.6pt
  \vskip8pt}

% (2) Ligatúry fi/fl sa pri kopírovaní a v ATS rozložia na bežné písmená
%     ("Artificial", nie "Arti<U+FB01>cial" -> ATS inak nenájde kľúčové slovo)
\ifpdftex
  \pdfglyphtounicode{f_i}{0066 0069}\pdfglyphtounicode{fi}{0066 0069}
  \pdfglyphtounicode{f_l}{0066 006C}\pdfglyphtounicode{fl}{0066 006C}
  \pdfglyphtounicode{f_f}{0066 0066}\pdfglyphtounicode{ff}{0066 0066}
  \pdfglyphtounicode{f_f_i}{0066 0066 0069}\pdfglyphtounicode{ffi}{0066 0066 0069}
  \pdfglyphtounicode{f_f_l}{0066 0066 006C}\pdfglyphtounicode{ffl}{0066 0066 006C}
  % Lato pomenúva ligatúry uniFB00–uniFB04
  \pdfglyphtounicode{uniFB00}{0066 0066}\pdfglyphtounicode{uniFB01}{0066 0069}
  \pdfglyphtounicode{uniFB02}{0066 006C}\pdfglyphtounicode{uniFB03}{0066 0066 0069}
  \pdfglyphtounicode{uniFB04}{0066 0066 006C}
\fi

% (3) Menej pretečených riadkov pri zarovnaní vľavo
\setlength{\emergencystretch}{1.5em}

%=== POMOCNÉ MAKRÁ ===========================================================
% Ikona, ktorá sa pri kopírovaní / v ATS neprepíše ako "\csuse {email symbol}:"
\newcommand{\icon}[1]{\BeginAccSupp{method=plain,ActualText={}}#1\EndAccSupp{}}

% Placeholder na doplnenie (v hotovom CV nesmie ostať!)
\newcommand{\fillin}[1]{{\color{Gold}\icon{\faPen}\,\textit{#1}}}

% Kontakt v hlavičke: \cvinfo[odkaz]{ikona}{text}
\newcommand{\cvinfo}[3][]{%
  \mbox{\textcolor{accent}{\icon{\normalfont#2}}~\ifstrempty{#1}{#3}{\href{#1}{#3}}}\hspace{1.3em}}

% Nadpis sekcie s ikonou + záložka (bookmark) v PDF
\newcounter{cvsec}
% (voliteľný 1. argument = medzera nad nadpisom; prvé nadpisy v stĺpcoch ju nemajú,
%  aby boli oba stĺpce presne v jednej rovine)
\newcommand{\cvsect}[3][\bigskip]{%
  \par\needspace{10\baselineskip}\nointerlineskip#1
  \stepcounter{cvsec}\pdfbookmark[0]{#3}{cvsec.\thecvsec}%
  {\color{heading}\cvsectionfont\icon{\makebox[1.5em][l]{\smash{\color{accent}\large#2}}}\MakeUppercase{#3}\par}%
  \vspace{-0.5ex}{\color{headingrule}\rule{\linewidth}{1.6pt}\par}\medskip}

% Pozícia: \cvrole{Pozícia}{Firma}{Obdobie}{Miesto}
\newcommand{\cvrole}[4]{%
  \par\needspace{7\baselineskip}\noindent{\large\bfseries\color{emphasis}#1}\hfill{\small\bfseries\color{body}#3}\par
  \vspace{2pt}\noindent{\bfseries\color{accent}#2}\hfill{\small\color{body}\icon{\faMapMarker*}~#4}\par
  \smallskip}

% Technológie pod položkou
\newcommand{\cvtech}[1]{\par\nopagebreak\vspace{1pt}\noindent{\small\color{body}\icon{\color{accent}\faCode}~\textit{#1}}\par}

% Projekt: \cvproject{Názov}{URL}{Rok}{Popis}{Technológie}
\newcommand{\cvproject}[5]{%
  \par\needspace{5\baselineskip}\noindent{\bfseries\color{emphasis}\href{#2}{#1~\icon{\color{accent}\faGithub}}}\hfill{\small\color{body}#3}\par
  \vspace{2pt}#4\cvtech{#5}}

% Ocenenie: \cvaward{ikona}{Čo}{Kde, rok}
\newcommand{\cvaward}[3]{%
  \par\noindent\makebox[1.8em][l]{\icon{\color{accent}#1}}%
  \parbox[t]{\dimexpr\linewidth-1.8em\relax}{\LaTeXraggedright\leavevmode{\bfseries\color{emphasis}#2}\par{\small #3}}\par\vspace{5pt}}

% Vzdelanie: \cvedu{Titul}{Škola}{Obdobie}{Detaily}
\newcommand{\cvedu}[4]{%
  \par\needspace{5\baselineskip}\noindent{\LaTeXraggedright\bfseries\color{emphasis}#1\par}
  \vspace{1pt}\noindent{\LaTeXraggedright\small\bfseries\color{accent}#2\par}
  \vspace{1pt}\noindent{\small\color{body}#3}\par
  \ifstrempty{#4}{}{\vspace{3pt}{\small #4}\par}\smallskip}

% Štítky: plný = hlavná zručnosť, obrys = ďalšie
\renewcommand{\cvtag}[1]{\tikz[baseline]\node[anchor=base,draw=body!30,fill=white,rounded corners=2pt,
  inner xsep=0.65ex,inner ysep=0.5ex,text height=1.5ex,text depth=.25ex,font=\small]{#1};}
\newcommand{\cvtagmain}[1]{\tikz[baseline]\node[anchor=base,draw=accent,fill=accent,text=white,rounded corners=2pt,
  inner xsep=0.65ex,inner ysep=0.5ex,text height=1.5ex,text depth=.25ex,font=\small\bfseries]{#1};}
\newcommand{\skillgroup}[1]{\par\medskip\noindent{\scriptsize\bfseries\color{accent}\textls[120]{\MakeUppercase{#1}}}\par\vspace{3pt}}
% ragged2e by inak nevedel zalomiť riadok medzi štítkami (pretekali za okraj)
\newenvironment{tags}{\par\LaTeXraggedright}{\par}

% Jazyk podľa CEFR: \cvlang{Jazyk}{úroveň}{počet dielikov: A1=1 ... C2/rodný=6}
\newlength{\cvlangseg}\newlength{\cvlanggap}\setlength{\cvlanggap}{2pt}
\newcommand{\cvlang}[3]{%
  \par\noindent\makebox[0.3\linewidth][l]{\bfseries\color{emphasis}#1}%
  \setlength{\cvlangseg}{\dimexpr(0.46\linewidth-5\cvlanggap)/6\relax}%
  \begin{tikzpicture}[baseline=-0.2ex]
    \foreach \i in {1,...,6}{%
      \ifnum\i>#3\relax\def\segcol{body!15}\else\def\segcol{accent}\fi
      \fill[\segcol] ({(\i-1)*(\cvlangseg+\cvlanggap)},0) rectangle ++(\cvlangseg,4pt);}
  \end{tikzpicture}\hfill{\small\color{body}#2}\par\vspace{4pt}}

%=== ČASOVÁ OS (TikZ) ========================================================
% POZOR: mesiace sú ORIENTAČNÉ – pôvodné CV obsahuje iba roky. Uprav podľa reality.
% Formát: rok + (mesiac-1)/12, napr. september 2021 = 2021.67
\newcommand{\tlnow}{2026.77}   % dnes (október 2026)
\newlength{\tlunit}
\newcommand{\tlbar}[6]{% {y}{od}{do}{text}{farba výplne}{farba textu}
  \fill[#5,rounded corners=2pt] ({#2-2021+0.012},{#1-0.16}) rectangle ({#3-2021-0.012},{#1+0.16});
  \node[anchor=west,inner sep=0pt,xshift=2.5pt,text=#6,font=\scriptsize\bfseries] at ({#2-2021},#1) {#4};}
\newcommand{\cvtimeline}{%
  \setlength{\tlunit}{\dimexpr(\linewidth-0.9cm)/6\relax}%
  \noindent\BeginAccSupp{method=escape,ActualText={Timeline of study and work 2021-2026 (details below)}}%
  \begin{tikzpicture}[x=\tlunit,y=1cm]
    \foreach \yr in {2021,...,2026}{%
      \draw[body!18] ({\yr-2021},0.36) -- ({\yr-2021},-1.44);
      \node[anchor=south west,inner sep=1pt,text=body,font=\scriptsize] at ({\yr-2021},0.36) {\yr};}
    \draw[body!18] (6,0.36) -- (6,-1.44);
    % dráhy: štúdium / výskum a pobyty / práca (2 riadky, lebo sa práce prekrývajú)
    \node[text=accent,font=\small] at (-0.16,0) {\icon{\faGraduationCap}};
    \node[text=accent,font=\small] at (-0.16,-0.42) {\icon{\faFlask}};
    \node[text=accent,font=\small] at (-0.16,-1.05) {\icon{\faBriefcase}};
    % štúdium
    \tlbar{0}{2021.67}{2024.50}{Bc. \textperiodcentered{} UPJŠ Košice}{GoldenEarth}{Ink}
    \tlbar{0}{2024.67}{2025.58}{Mgr. \textperiodcentered{} Charles Univ.}{GoldenEarth}{Ink}
    \tlbar{0}{2025.67}{\tlnow}{PhD \textperiodcentered{} UPJŠ}{GoldenEarth}{Ink}
    % študentská vedecká činnosť a Erasmus (letný semester 2024 – odhad)
    \tlbar{-0.42}{2022.25}{2023.40}{Student research (ŠVOČ)}{GoldenEarth!70}{Ink}
    \tlbar{-0.42}{2024.08}{2024.55}{Erasmus+}{GoldenEarth!70}{Ink}
    % práca – riadok 1
    \tlbar{-0.84}{2021.50}{2024.60}{KDC \textperiodcentered{} C\#/.NET, JavaScript}{accent!75}{white}
    \tlbar{-0.84}{2025.00}{2025.45}{DADO}{accent!75}{white}
    \tlbar{-0.84}{2025.50}{\tlnow}{LevelGo}{accent}{white}
    % práca – riadok 2
    \tlbar{-1.26}{2026.00}{2026.42}{Instaview}{accent!75}{white}
    \tlbar{-1.26}{2026.44}{\tlnow}{UNLP}{accent}{white}
    % dnes
    \draw[accent,densely dashed,thick] ({\tlnow-2021},0.28) -- ({\tlnow-2021},-1.42);
    \node[anchor=south,inner sep=1pt,text=accent,font=\scriptsize\itshape] at ({\tlnow-2021},0.36) {today};
  \end{tikzpicture}\EndAccSupp{}\par}

%=== PUBLIKÁCIE ===============================================================
\usepackage[backend=biber,style=ieee,sorting=ydnt,defernumbers=true]{biblatex}
%% For removing numbering entirely when using a numeric style
\setlength{\bibhang}{1.25em}
\DeclareFieldFormat{labelnumberwidth}{\makebox[\bibhang][l]{\itemmarker}}
\setlength{\biblabelsep}{0pt}
\defbibheading{pubtype}{\cvsubsection{#1}}
\renewcommand{\bibsetup}{\vspace*{-\baselineskip}}
\AtEveryBibitem{%
  \iffieldundef{doi}{}{\clearfield{url}}%
}

\addbibresource{publications.bib}

\begin{document}

%=== HLAVIČKA ================================================================
\noindent
\begin{minipage}[c]{\dimexpr\linewidth-2.4cm\relax}
  {\namefont\color{name}NORBERT MICHEĽ\par}
  \smallskip
  {\taglinefont\color{tagline}Software Developer \& Researcher \enspace\textperiodcentered\enspace LLM-Powered Algorithms\par}
  \medskip
  {\footnotesize\color{body}%
    \cvinfo[mailto:noro.michel159@gmail.com]{\faEnvelope}{noro.michel159@gmail.com}%
    \cvinfo[tel:+421907869883]{\faPhone}{+421 907 869 883}%
    \cvinfo{\faMapMarker*}{Košice, Slovakia}%
    \cvinfo[https://www.linkedin.com/in/norbert-miche\%C4\%BE]{\faLinkedin}{norbert-micheľ}%
    \cvinfo[https://github.com/nmichel159]{\faGithub}{nmichel159}\par}
\end{minipage}\hfill
\begin{minipage}[c]{2.1cm}\centering
  \qrcode[height=1.9cm]{https://github.com/nmichel159}\par\vspace{2pt}
  {\tiny\color{body}GitHub}
\end{minipage}\par

\medskip
\cvtimeline
\vspace{9pt}

%=== DVA STĹPCE ================================================================
\columnratio{0.62}
\backgroundcolor{c[1](7pt,2pt)(1.25cm,1.0cm)}{Sidebar}
\begin{paracol}{2}

\cvsect[]{\faBriefcase}{Experience}

\cvrole{Optimization Engineer}{UNLP -- Louis Pasteur University Hospital}{2026 -- present}{Košice}
\begin{itemize}
  \item Building a shift-scheduling system for hospital staff that generates duty rosters using integer linear programming (ILP).
  \item Developing the full application around the optimization model -- Python back end (FastAPI) and React front end -- with an agentic AI-driven workflow.
\end{itemize}
\cvtech{Python \textperiodcentered{} FastAPI \textperiodcentered{} React \textperiodcentered{} integer linear programming\enspace
  \href{https://github.com/nmichel159/scheduling}{\icon{\faGithub}}}
\divider

\cvrole{Software Developer}{LevelGo}{2025 -- present}{Košice}
\begin{itemize}
  \item Developing a Flutter sports app built on clean architecture.
  \item Building the Python back-end services: application logic, data processing and the API used by the mobile client.
\end{itemize}
\cvtech{Flutter (Dart) \textperiodcentered{} Python}
\divider

\cvrole{AI Engineer}{Instaview}{2026}{Prague}
\begin{itemize}
  \item Developed an LLM-based system that analyzes job interviews and evaluates candidates' answers for both technical and soft skills.
  \item Designed algorithms that match candidates to open positions based on skill alignment and the semantic meaning of their answers.
\end{itemize}
\cvtech{Python \textperiodcentered{} LLMs}
\divider

\cvrole{Full-Stack Developer}{DADO}{2025}{Bratislava}
\begin{itemize}
  \item Integrated third-party AI APIs to automate data processing and content generation.
  \item Tailored AI tools to client requirements and workflow automation.
\end{itemize}
\cvtech{Vue.js \textperiodcentered{} PHP \textperiodcentered{} AI APIs}
\divider

\cvrole{Software Developer}{KDC -- Košice Development Center}{2021 -- 2024}{Košice}
\begin{itemize}
  \item Built a full-stack C\# application for Slovak Railways (ZSSK), including data synchronization with ZSSK's internal system.
  \item Designed a reusable Blazor component library (counters, sliders and other UI controls) with demo apps for internal developers.
  \item Built and maintained front ends and API integrations for several government-related websites.
\end{itemize}
\cvtech{C\# \textperiodcentered{} .NET \textperiodcentered{} Blazor \textperiodcentered{} JavaScript}

\cvsect{\faLaptopCode}{Selected Projects}

\cvproject{Skill Analyzer (LLM)}{https://github.com/nmichel159/skill-analyzer-LLM}{2026}
  {Uses large language models to build a skill-similarity matrix that helps database queries compare and match skills.}
  {Python \textperiodcentered{} LLMs \textperiodcentered{} PostgreSQL}
\divider
\cvproject{IYPT Tournament Scheduler}{https://github.com/nmichel159/IYPT_Linear_programming}{2022}
  {Student research project: randomized algorithms combined with linear programming make an NP-hard scheduling problem tractable in practice -- fair, efficient match schedules for the International Young Physicists' Tournament.}
  {Python \textperiodcentered{} PuLP \textperiodcentered{} Tkinter}
\divider
\cvproject{3D Tic-Tac-Toe Dropped}{https://github.com/nmichel159/3D_tic_tac_toe_dropped}{2022}
  {3D ``falling'' tic-tac-toe (3$\times$3$\times$3 up to 5$\times$5$\times$5, gravity as in Connect Four) with custom rendering: a vanishing-point projection and its inverse map the 3D board onto the 2D canvas and back for click detection.}
  {Python \textperiodcentered{} Tkinter \textperiodcentered{} projective geometry}
\divider
\cvproject{Distressed-Value Stock Screener}{https://github.com/nmichel159/stock}{2026}
  {Streamlit research tool for finding financially distressed but potentially undervalued stocks across global exchanges.}
  {Python \textperiodcentered{} Streamlit \textperiodcentered{} Pandas}
\divider
\cvproject{Sport App}{https://github.com/nmichel159/sport}{2026}
  {Sports activity management app that works both offline and online: authentication, match management and statistics.}
  {TypeScript \textperiodcentered{} React Native \textperiodcentered{} FastAPI}

\cvsect{\faBookOpen}{Publications}

\mynames{Micheľ/N.}
\nocite{*}
\printbibliography[heading=pubtype,title={Manuscripts in preparation},type=unpublished]

\cvsect{\faChalkboardTeacher}{Mentoring}

\cvedu{Construction of a Small-Cap Equity Investment Portfolio Using Large Language Models}{Master's thesis -- Kiara Koščová}{2026}{}
\divider
\cvedu{Algorithms for Move Generation and Evaluation in Chess}{Bachelor's thesis -- Tomáš Oltman}{2025}{}

%=== PRAVÝ STĹPEC ==============================================================
\switchcolumn

\cvsect[]{\faTrophy}{Highlights}

\cvaward{\faBolt}{Mgr. in 1 year (standard 2)}{Charles University, Prague}
\cvaward{\faMedal}{Silver Medal -- MEMO 2020}{Middle European Mathematical Olympiad}
\cvaward{\faTrophy}{Winner -- National Olympiads}{Mathematical Olympiad \& Olympiad in Informatics}
\cvaward{\faLaptopCode}{Slovak Champion -- CERC 2022}{ICPC Central Europe Regional Contest}
\cvaward{\faAward}{Bernard Bolzano Prize nominee 2025}{Best-thesis award, Charles University (MFF)}
\cvaward{\faCode}{2nd Place -- UPJŠ Contest 2022}{Student programming competition, Košice}

\cvsect{\faGraduationCap}{Education}

\cvedu{PhD in Computer Science}{P. J. Šafárik University in Košice}{2025 -- present}
  {Thesis: \emph{Optimisation Algorithms Powered by LLMs}\\ Supervisor: prof.\ RNDr.\ Gabriel Semanišin, PhD.}
\divider
\cvedu{Mgr. (MSc) in Discrete Models and Algorithms}{Charles University, Faculty of Mathematics and Physics, Prague}{2024 -- 2025 \textperiodcentered{} 1 year instead of 2}
  {Thesis: \emph{Min Cut-Path} -- nominated for the Bernard Bolzano Foundation Prize\\ Supervisor: prof.\ RNDr.\ Martin Loebl, CSc.}
\divider
\cvedu{Bc. in Data Analysis and Artificial Intelligence}{P. J. Šafárik University in Košice}{2021 -- 2024}
  {Thesis: \emph{Lattice Polytopes in the $n$-Dimensional Cube}\\ Supervisor: Mgr.\ Martin Vodička}

\cvsect{\faLanguage}{Languages}

\cvlang{Slovak}{native}{6}
\cvlang{Czech}{C1}{5}
\cvlang{English}{B2}{4}
\cvlang{German}{A2}{2}

\cvsect{\faTools}{Skills}

\skillgroup{Agentic AI development}
\noindent\makebox[1.8cm][l]{\LARGE\bfseries\color{accent}15B+}%
\parbox[c]{\dimexpr\linewidth-1.8cm\relax}{\small tokens processed with AI coding agents}\par\vspace{5pt}
\begin{tags}
\cvtagmain{Claude Code}
\cvtag{Codex}
\cvtag{Antigravity}
\cvtag{Agentic workflows}
\end{tags}

\skillgroup{Programming languages}
\begin{tags}
\cvtagmain{Python}
\cvtagmain{C++}
\cvtagmain{C\#}
\cvtagmain{TypeScript}
\cvtag{JavaScript}
\cvtag{Java}
\cvtag{Kotlin}
\cvtag{Dart}
\cvtag{Rust}
\cvtag{Haskell}
\cvtag{C}
\cvtag{PHP}
\cvtag{SQL}
\cvtag{Bash}
\end{tags}

\skillgroup{AI \& LLMs}
\begin{tags}
\cvtagmain{LLMs}
\cvtagmain{Prompt engineering}
\cvtagmain{RAG}
\cvtag{Embeddings \& semantic search}
\cvtag{OpenAI API}
\cvtag{LangChain}
\cvtag{Hugging Face}
\cvtag{Ollama}
\cvtag{PyTorch}
\cvtag{scikit-learn}
\cvtag{NLP}
\end{tags}

\skillgroup{Data science}
\begin{tags}
\cvtagmain{NumPy}
\cvtagmain{Pandas}
\cvtag{SciPy}
\cvtag{Matplotlib}
\cvtag{NetworkX}
\cvtag{Jupyter}
\cvtag{Streamlit}
\end{tags}

\skillgroup{Optimization \& algorithms}
\begin{tags}
\cvtagmain{Linear \& integer programming}
\cvtagmain{Gurobi}
\cvtag{PuLP}
\cvtagmain{Combinatorial optimization}
\cvtagmain{Graph algorithms}
\cvtag{Randomized algorithms}
\cvtag{Algorithms \& data structures}
\end{tags}

\skillgroup{Web \& mobile}
\begin{tags}
\cvtagmain{FastAPI}
\cvtag{Flask}
\cvtag{Node.js}
\cvtag{ASP.NET Core}
\cvtag{Blazor}
\cvtag{REST APIs}
\cvtagmain{React}
\cvtag{Vue}
\cvtag{Angular}
\cvtag{HTML / CSS}
\cvtagmain{Flutter}
\cvtag{React Native}
\cvtag{Expo}
\cvtag{Android}
\end{tags}

\skillgroup{Databases}
\begin{tags}
\cvtagmain{PostgreSQL}
\cvtag{MySQL}
\cvtag{SQL Server}
\cvtag{SQLite}
\end{tags}

\skillgroup{DevOps \& tools}
\begin{tags}
\cvtagmain{Git / GitHub}
\cvtagmain{Docker}
\cvtag{Linux}
\cvtag{AWS}
\cvtag{Railway}
\cvtag{LaTeX}
\end{tags}

\cvsect{\faHeart}{Interests}

\begin{tags}
\cvtag{Chess (FIDE 2000+)}
\cvtag{Ultimate Frisbee}
\cvtag{Bouldering}
\end{tags}

% bočný panel dotiahne farbu až po spodok poslednej strany
\par\vspace*{\dimexpr\pagegoal-\pagetotal-2pt\relax}

\end{paracol}
\end{document}
